\chapter{Обзор предметной области}\label{ch:review}

\section{Архитектура электрофоретического дисплея}\label{sec:epd-arch}

Электрофоретический дисплей (Electrophoretic Display, EPD) --- отражательное
устройство визуализации, изображение в котором формируется позиционированием
заряженных пигментных частиц внутри микрокапсул, заполненных непроводящей
жидкостью. Технология введена Comiskey~и~соавторами в 1998~году~\cite{comiskey1998_nature},
впервые продемонстрировавшими изготовление отражательного дисплея методами
печатной технологии. С тех пор EPD заняли устойчивую нишу в продуктах,
требующих низкого среднего энергопотребления и читаемости при внешнем
освещении: ридерах электронных книг (Amazon~Kindle, Kobo, ONYX~BOOX),
электронных ценниках, низкоэнергетических индикаторах для интернета вещей и
носимой электроники.

Принцип работы микрокапсулы основан на электрофоретическом смещении двух типов
частиц с противоположными знаками заряда (обычно белых частиц диоксида
титана~$\mathrm{TiO}_2$ и чёрных частиц технического углерода) под действием
поля на электродах пикселя~\cite{yang2021_mechanisms}. В активных матрицах
(active matrix EPD, AM\nobreakdash-EPD) индивидуальное управление нижними
электродами реализуется тонкоплёночными транзисторами; общая структура
устройства показана на рисунке~\ref{fig:02_am_epd_arch}.

\begin{figure}[H]
    \centering
    \begin{tikzpicture}[font=\small, scale=1.0]
        % Верхний прозрачный ITO электрод
        \draw[fill=cyan!10, draw=cyan!60!black, thick]
            (0, 2) rectangle (9, 2.4);
        \node[right=0.1cm] at (9, 2.2) {ITO common electrode};

        % Слой микрокапсул (3 капсулы)
        \foreach \x in {1.5, 4.5, 7.5} {
            \draw[fill=gray!5, draw=gray!50, thick] (\x, 1) circle (1.0);
        }

        % Частицы внутри: белые сверху, чёрные снизу
        \foreach \x in {1.5, 4.5, 7.5} {
            % белые частицы (T iO2) — наверху
            \foreach \y/\xx in {1.6/0, 1.7/0.3, 1.5/-0.4, 1.6/-0.2} {
                \draw[fill=white, draw=black, thin] (\x+\xx, \y) circle (0.10);
            }
            % чёрные частицы — внизу
            \foreach \y/\xx in {0.6/0.1, 0.55/-0.3, 0.65/-0.05, 0.5/0.35} {
                \draw[fill=black] (\x+\xx, \y) circle (0.10);
            }
        }

        % Нижний электрод (TFT матрица — сегментированный)
        \foreach \x in {0, 1.5, 3, 4.5, 6, 7.5} {
            \draw[fill=yellow!20, draw=brown, thick]
                (\x+0.1, -0.3) rectangle (\x+1.4, 0);
        }
        \node[below=0.1cm] at (4.5, -0.3) {TFT-подложка (сегментированная)};

        % Управляющее напряжение
        \draw[<->, thick, gray] (-0.5, 0) -- (-0.5, 2.4)
            node[midway, left] {\rotatebox{90}{$U(t)$}};

        % Метки сторон
        \node[above right] at (0, 2.4) {\small (+)};
        \node[below right] at (0, -0.3) {\small (--)};
    \end{tikzpicture}
    \caption{Архитектура активно-матричного электрофоретического дисплея.
        Сверху --- общий прозрачный ITO электрод, снизу --- сегментированная
        TFT-матрица (подложка). Между электродами --- слой микрокапсул с
        белыми (TiO$_2$) и чёрными частицами противоположных зарядов.
        Управляющее напряжение $U(t)$ задаёт направление электрофоретического
        смещения частиц и тем самым видимое оптическое состояние пикселя.}
    \label{fig:02_am_epd_arch}
\end{figure}
\FloatBarrier

Ключевое физическое свойство EPD --- \emph{бистабильность}: после снятия
управляющего поля частицы сохраняют положение часами и днями без
энергопотребления~\cite{yang2021_mechanisms}. Это отличает EPD от
жидкокристаллических и органических светодиодных дисплеев, где поддержание
изображения требует непрерывного тока. За счёт бистабильности в статическом режиме
энергопотребление близко к нулю, а заметные затраты приходятся только на момент
обновления.

С точки зрения моделирования EPD относятся к классу
электрофоретико-диэлектрофоретических систем. Помимо силы
электрофореза~$\vec{F}_{\!\text{EP}} = q \vec{E}$ на заряженную частицу, в
неоднородном поле существенна сила диэлектрофореза
$\vec{F}_{\!\text{DEP}} \propto \nabla |\vec{E}|^2$, объясняющая отклик пикселей
на переменное поле и эффективный порог переключения при пассивной
адресации~\cite{bert2003_physics}. В рассматриваемом здесь режиме
(последовательность постоянных импульсов с обнулённым средним зарядом) вкладом
DEP можно пренебречь, что обсуждается в главе~\ref{ch:model}.


\section{Таксономия методов проектирования waveform}\label{sec:taxonomy}

Под \emph{waveform} в EPD понимается последовательность дискретных фаз
$(V_k, T_k)_{k=1}^{K}$, прикладываемых к пикселю при переходе между уровнями
серого. Её конфигурация --- основной (и обычно единственный программно
доступный) инструмент оптимизации обновления EPD на уровне контроллера.
Накопленный за два десятилетия пул работ систематизирован в
обзоре~\cite{zhong2026_review}; на его основе ниже выделены пять
направлений, наиболее значимых для постановки настоящей работы.

\subsection*{Исторический фундамент}

Базовая трёхстадийная структура waveform (\emph{erase} — \emph{activate} —
\emph{drive}) предложена Zehner~и~соавторами в 2003~году (цитируется
по~\cite{kang2025_dp_waveform}). Все последующие модификации waveform для
b/w-панелей сохраняют эту трёхчастную базу, варьируя структуру отдельных стадий.

\subsection*{Подавление ghosting}

Wang~и~соавторы~\cite{wang2022_red_ghost} снизили «красный гостинг» в
трёхцветных EPD, разделив этап стирания на две части (специализированную для
красных частиц и стандартную) и применив высокочастотный прямоугольный сигнал
на этапе активации. Подавление остаточного изображения составило 80{,}43\,\%,
интенсивности фликера — 79{,}63\,\%. Авторы отдельно оговаривают
требование нулевого суммарного заряда (DC balance compliance),
вытекающее из физики накопления заряда в микрокапсуле и подтверждаемое
сокращением ресурса панели при несбалансированных
waveform~\cite{yang2021_mechanisms}.

\subsection*{Ускорение отклика}

He~и~соавторы~\cite{he2020_particle_activation} разделили активационную фазу на
две подфазы: повышение подвижности частиц и установление референсного уровня
серого. Длительность первой подфазы выбирается по точке перегиба эмпирической
кривой яркости (момент максимума скорости частиц). На панели~ED060SC7 это
сократило полное время обновления на~300\,мс при подавлении гостинга на~57\,\%
и снижении фликера на~26{,}7\,\%.

\subsection*{Снижение энергопотребления}

Наиболее свежие работы по энергоэффективности --- Lin~и~соавторов~\cite{lin2024_low_power}
и Lai~и~соавторов~\cite{lai2026_e3dw_energy}.

Lin~и~соавторы исходят из соотношения для пиковой мощности обновления:
\begin{equation}\label{eq:p-peak}
    P_{\text{peak}} \;=\; \tfrac{1}{2}\,C\,\Delta V^{2}\, f\, V_{\text{source}},
\end{equation}
где $C$~--- эквивалентная ёмкость пикселя, $\Delta V$~--- амплитуда скачка
напряжения между соседними фазами, $f$~--- частота переключений,
$V_{\text{source}}$~--- максимальное напряжение источника.\footnote{Соотношение
приводится в записи первоисточника~\cite{lin2024_low_power}; размерно оно
отвечает мощности с точностью до множителя $V_{\text{source}}$ и используется
авторами как качественный индикатор зависимости потребления от $\Delta V$ и
$f$. В собственной энергетической модели настоящей работы (гл.~\ref{ch:experiment})
применяется калибровка на уровне панели, а не эта формула.} Уменьшение числа
смен полярности и нулевые вставки в waveform снизили энергетическую флуктуацию
на~37{,}24\,\% и среднюю мощность на тестовом наборе из девяти изображений.

Lai~и~соавторы предлагают системно-уровневое решение E3DW
(\emph{Energy-Efficient E-Paper Driving Waveform}), объединяющее анализ
содержимого кадра, компенсацию по температуре и низкочастотный прямоугольный
сигнал (1--5\,Гц). На цветном EPD достигнуто снижение энергопотребления
на~35{,}7\,\% относительно базовой конфигурации.

\subsection*{Алгоритмическая оптимизация на основе динамического программирования}

Kang~и~соавторы~\cite{kang2025_dp_waveform} представляют наиболее близкий к
настоящей работе подход: строят дискретный граф состояний пикселя и применяют
динамическое программирование для поиска оптимального пути по скаляризованному
критерию --- взвешенной сумме времени отклика, фликера и качества изображения
с весами~$0{,}6$, $0{,}2$ и~$0{,}2$. Главное ограничение (см.~раздел~\ref{sec:gap}) --- фиксация весов: решение есть единственная точка в
пространстве~$(E, G, \tau)$, без исследования компромиссов между целями.

\subsection*{Сравнительная таблица}

В таблице~\ref{tab:waveform-methods} собраны рассмотренные методы с указанием
оптимизируемых величин, алгоритмической основы и заявленных результатов;
в последней строке для сопоставления приведена позиция настоящей работы.

\begin{table}[H]
    \caption{Сравнение современных методов проектирования waveform для EPD}\label{tab:waveform-methods}
    \begin{tabular}{l p{3.5cm} p{3.5cm} p{3.5cm}}
        \toprule
        Источник & Что оптимизируется & Алгоритмическая основа & Результат \\
        \midrule
        Zehner~2003 (cit.\ \cite{kang2025_dp_waveform}) & базовая структура waveform & ручной дизайн & базовая трёхстадийная схема \\
        He~2020~\cite{he2020_particle_activation} & время отклика & эвристика точки перегиба & $-300$\,мс, $-57\,\%$ ghost \\
        Wang~2022~\cite{wang2022_red_ghost} & красный гостинг (3-color) & разделение стадии стирания & $-80{,}4\,\%$ ghost \\
        Lin~2024~\cite{lin2024_low_power} & пиковая мощность & вставки нулевого напряжения & $-37\,\%$ флуктуация энергии \\
        Lai~2026~\cite{lai2026_e3dw_energy} & энерго\-потребление & system-level E3DW & $-35{,}7\,\%$ энергии \\
        Kang~2025~\cite{kang2025_dp_waveform} & скаляризованный $(\tau, G, \text{кач.})$ & динамическое программирование & оптимальный путь, single objective \\
        \textbf{Настоящая работа} & \textbf{Парето-граница $(E, G, \tau)$} & \textbf{$\varepsilon$-constraint + дискретный PMP} & \textbf{весь Парето-фронт, а не одна точка} \\
        \bottomrule
    \end{tabular}
\end{table}


\section{Контроллер SSD1680}\label{sec:ssd1680}

Аппаратной базой работы служит контроллер SSD1680 производства
Solomon~Systech~\cite{ssd1680_datasheet}, управляющий панелью
Waveshare~2.13''~V4. Его архитектура задаёт конечномерную дискретную постановку
задачи оптимального управления, формализуемую в~главе~\ref{ch:theorem}.

\subsection*{Структура waveform setting и регистр 0x32}

Согласно разделу~6.7 datasheet, полный waveform setting состоит из~159 байт,
кодирующих все управляющие параметры одного обновления: 153~байта LUT
(записываются командой~$0x32$) и по одному байту для выходных напряжений
$V_{\text{GH}}$, $V_{\text{SH1}}$, $V_{\text{SH2}}$, $V_{\text{SL}}$,
$V_{\text{COM}}$ плюс опциональный байт окончания процедуры. Структура LUT ---
5 sub-LUT $\times$ 12 фаз $\times$ 4 sub-frame; в каждой ячейке выбирается один
из пяти уровней напряжения
$\{V_{\text{SH1}}, V_{\text{SH2}}, V_{\text{SL}}, V_{\text{COM}}, 0\}$, что
задаёт конечный алфавит управляющих воздействий мощностью~$|U| = 5$.

Для каждой фазы дополнительно указываются длительности обоих sub-frame'ов
$\mathrm{TP}_k^{X}$, параметры наклона переходов~$\mathrm{SR}_k^{XY}$, число
повторений фазы $\mathrm{RP}_k$, частота кадров $\mathrm{FR}_k$ и битовый флаг
gate-gating $\mathrm{XON}_k^{XY}$. В совокупности они задают конечномерное
пространство допустимых waveform с явной структурой управления.

\subsection*{Поиск waveform по температуре в OTP}

Помимо программно-записываемого LUT контроллер хранит во встроенной однократно
программируемой памяти (One-Time Programmable, OTP) до~36 температурно-индексированных
waveform setting (TR0\ldots TR35). Команды $0x18$ (выбор источника датчика
температуры) и $0x22$ с битом «загрузить значение температуры» инициируют поиск
подходящей записи в OTP и её загрузку в регистр~$0x32$~\cite{ssd1680_datasheet}.

\subsection*{Семантика регистра 0x22 и условие применимости пользовательской LUT}

Регистр $0x22$ (\emph{Display Update Control~2}) задаёт последовательность
операций при активации обновления командой~$0x20$. Восемь основных значений и
соответствующие последовательности приведены в
datasheet~\cite{ssd1680_datasheet}; для нас важны четыре варианта:

\begin{itemize}
    \item $0x{}C7$~/~$0x{}CF$ — последовательность \emph{«включить аналоговое
    питание $\to$ запустить обновление в режиме Display Mode 1/2 $\to$ выключить
    аналоговое питание $\to$ выключить генератор»}, использующая waveform,
    \emph{уже записанный пользователем} в регистр~$0x32$;
    \item $0x{}F7$~/~$0x{}FF$ — тот же набор операций с дополнительным шагом
    \emph{«загрузить значение температуры»}, инициирующим Waveform Setting
    Searching Mechanism и перезапись пользовательского LUT заводским вариантом
    из OTP.
\end{itemize}

Отсюда следует, что применение пользовательской waveform-LUT возможно
\emph{исключительно} через значение $0x{}C7$ (полное обновление) или $0x{}CF$
(частичное); использование $0x{}F7$~/~$0x{}FF$ эквивалентно отказу от
пользовательской waveform. Это наблюдение, отсутствующее в reference-реализации
Waveshare~\cite{waveshareteam_epaper}, воспроизведено и подтверждено в
экспериментах работы (см.~раздел~\ref{ch:experiment}). Без указанного фикса
любые алгоритмические улучшения waveform неприменимы на физическом стенде.


\section{Постановка задачи и её отличие от существующих работ}\label{sec:gap}

В рассмотренных в~разделе~\ref{sec:taxonomy} работах отсутствует ряд элементов,
существенных для систематической инженерной оптимизации обновления EPD.

Во-первых, ни в одной из них нет структурного результата о форме
оптимального waveform. Все цитированные методы предлагают алгоритмические
процедуры поиска: ручную настройку~\cite{wang2022_red_ghost}, эвристику точки
перегиба~\cite{he2020_particle_activation}, перебор
конфигураций~\cite{lin2024_low_power} или динамическое
программирование~\cite{kang2025_dp_waveform}. Знание структуры оптимума
(например, что оптимальное управление имеет релейную (bang-bang) структуру с
числом активных фаз, ограниченным размерностью пространства состояний)
позволило бы a~priori сузить пространство поиска и обосновать тип кандидатных
waveform.

Во-вторых, все эти работы оптимизируют одну целевую функцию или
фиксированную скаляризацию нескольких целей: Kang~и~соавторы~\cite{kang2025_dp_waveform}
--- веса $(0{,}6;\, 0{,}2;\, 0{,}2)$; \cite{lin2024_low_power} --- только
пиковую мощность; \cite{lai2026_e3dw_energy} --- суммарную энергию. При этом
авторы обзора~\cite{zhong2026_review} отмечают, что в реальных приложениях вес
каждой цели сильно зависит от сценария: для мобильных устройств критична
энергия, для считывающих --- латентность, для систем длительного отображения
--- гостинг. Единственной точки в пространстве $(E, G, \tau)$ недостаточно для
осмысленного адаптивного выбора waveform.

В-третьих, управляющие ограничения нигде не выводятся из
физико-математической модели устройства. Прежде всего это касается условия
нулевого суммарного заряда, \emph{зарядового баланса}. Оно упоминается как общепринятая
практика~\cite{wang2022_red_ghost, yang2021_mechanisms}, но не интегрируется в
формальную постановку оптимизационной задачи.

Наконец, валидация в~\cite{kang2025_dp_waveform,
lai2026_e3dw_energy} проводилась на специализированном оборудовании
(колориметры, лабораторные анализаторы мощности), что затрудняет независимое
воспроизведение и сопоставление методик.

Перечисленные пробелы и закрывает настоящая работа: каскадной моделью
обновления (глава~\ref{ch:model}), структурной теоремой о форме оптимума в классе
сбалансированных по заряду waveform (глава~\ref{ch:theorem}), алгоритмом
построения Парето-границы по $(E, G, \tau)$ на базе $\varepsilon$-constraint
метода (глава~\ref{ch:algorithm}) и экспериментальной валидацией на стенде из
общедоступных компонентов с открытым ПО (глава~\ref{ch:experiment}).
